\subsection{Problem Formulation and Inductive Evaluation Protocol}
\label{sec:problem}

Let $\mathcal{D} = \{(\bx_i, b_i)\}_{i=1}^N$ denote a single-cell dataset of $N$ cells, where $\bx_i \in \R^G$ is the log-normalized expression vector of cell $i$ and $b_i \in \{1, \ldots, B\}$ is its batch or donor label. A subset $\mathcal{L} \subset [N]$ of cells has known cell-type labels $y_i \in \calC = \{c_1, \ldots, c_K\}$. One class $\rare \in \calC$ is the \emph{rare class} of interest; all others are \emph{majority classes}. The training set $\calT$ consists of all labeled cells from a reference subset of batches; the remaining batches form the \emph{query set} $\mathcal{Q}$. Under the \emph{batch-held-out} evaluation protocol, no cell from any query batch appears in $\calT$.

Let $\bz_i = f_\phi(\bx_i, b_i) \in \R^d$ denote the $d$-dimensional latent representation produced by a pre-trained scANVI encoder with parameters $\phi$. scANVI also produces a softmax prediction $\hat{y}_i = \argmax_{c} P_\phi(c \mid \bz_i)$. Our goal is to produce a refined prediction $\tilde{y}_i$ that corrects cases where $\hat{y}_i \neq \rare$ but $y_i = \rare$.

\textbf{Inductive constraint}: All reference quantities used by scRareRefine---class prototypes, marker signatures, gate thresholds, and fusion parameters---are derived exclusively from $\calT$. Test-set labels are used only for final evaluation and never influence any algorithmic decision.

\subsection{Separability Ratio}
\label{sec:sepratio}

The \emph{prototype} of class $c$ is the centroid of its labeled training embeddings:
\begin{equation}
    \proto{c} = \frac{1}{|\calT_c|} \sum_{i \in \calT_c} \bz_i,
    \label{eq:prototype}
\end{equation}
where $\calT_c = \{i \in \calT : y_i = c\}$.

The \emph{intra-class radius} of the rare class measures cluster compactness in the training set:
\begin{equation}
    \intrarad = \frac{1}{|\calT_\rare|} \sum_{i \in \calT_\rare} \|\bz_i - \proto{\rare}\|_2.
    \label{eq:intrarad}
\end{equation}

The \emph{nearest-majority distance} is the Euclidean distance from the rare prototype to the closest majority-class prototype:
\begin{equation}
    \majdist = \min_{c \in \calC \setminus \{\rare\}} \|\proto{\rare} - \proto{c}\|_2.
    \label{eq:majdist}
\end{equation}

The \emph{separability ratio} is defined as:
\begin{equation}
    \sepratio = \frac{\majdist}{\intrarad}.
    \label{eq:sepratio}
\end{equation}

$\sepratio > 1$ indicates that the rare prototype is farther from the nearest majority prototype than the average within-cluster spread of rare training cells---a necessary geometric condition for prototype-based discrimination. We observe empirically that $\sepratio > 1.3$ reliably predicts positive F1 gain across all evaluated datasets (see \cref{sec:sepratio_result}). The separability ratio is computable entirely from $\calT$ before any query inference, making it a practical pre-deployment diagnostic.

\paragraph{Deployment thresholds.} We use two threshold levels:
\begin{itemize}[leftmargin=*]
    \item $\sepratio \geq 1.3$: the full rescue pipeline is activated.
    \item $\sepratio < 1.1$: the candidate set is emptied and all query cells retain their scANVI predictions (fallback mode). Intermediate values ($1.1 \leq \sepratio < 1.3$) proceed through the pipeline with conservative gate settings.
\end{itemize}

\subsection{Prototype Scoring and Candidate Selection}
\label{sec:scoring}

For each query cell $i$ with $\hat{y}_i \neq \rare$, we compute the Euclidean distance from $\bz_i$ to each class prototype and rank the rare class:
\begin{equation}
    \text{rank}_i = \text{rank}_\rare\!\left(\{\|\bz_i - \proto{c}\|_2\}_{c \in \calC}\right),
    \label{eq:rank}
\end{equation}
where rank 1 indicates that the rare prototype is the nearest centroid. We also compute the \emph{distance margin}:
\begin{equation}
    m_i = \|\bz_i - \proto{\hat{y}_i}\|_2 - \|\bz_i - \proto{\rare}\|_2,
    \label{eq:margin}
\end{equation}
measuring how much closer the rare prototype is compared with the predicted-class prototype (positive values indicate the rare prototype is closer).

A query cell is designated a \emph{rescue candidate} if both of the following conditions hold:
\begin{enumerate}[leftmargin=*]
    \item $\text{rank}_i \leq 2$ (the rare prototype is the first or second nearest centroid); and
    \item $m_i > q_{25}(m_{\calT})$, where $q_{25}(m_{\calT})$ is the 25th percentile of the margin distribution computed over cells in $\calT$ (using leave-one-out distances for labeled cells).
\end{enumerate}
The margin threshold is derived from the training set and is not tuned on validation data.

\subsection{Prototype Gate and Fallback Mechanism}
\label{sec:gate}

Before proceeding to marker verification, a structural gate checks the separability ratio:
\begin{itemize}[leftmargin=*]
    \item If $\sepratio < 1.1$: the candidate set is set to empty, and all cells retain their scANVI predictions.
    \item If $1.1 \leq \sepratio < 1.3$: candidates proceed but with an additional constraint that the distance margin $m_i$ must exceed the 50th (rather than 25th) percentile threshold.
    \item If $\sepratio \geq 1.3$: candidates proceed with the standard margin threshold as described above.
\end{itemize}

This hard gate prevents over-rescue in datasets where the latent-space geometry does not support prototype discrimination, which is the primary cause of false-rescue inflation in purely distance-based methods. It also ensures that the method produces no net harm when the rare class is not geometrically separable.

\subsection{Marker-Gene Verification}
\label{sec:marker}

For each rescue candidate, we apply a biological plausibility check using a differential expression signature derived from $\calT$. For each class $c$ with $|\calT_c| \geq 5$ labeled training cells, we select the top-$k$ discriminative genes:
\begin{equation}
    \calG_c = \underset{g : \mu_{c,g} > \bar{\mu}_{-c,g}}{\text{top-}k} \left(\mu_{c,g} - \bar{\mu}_{-c,g}\right),
    \label{eq:markers}
\end{equation}
where $\mu_{c,g} = \frac{1}{|\calT_c|}\sum_{i \in \calT_c} x_{ig}$ is the mean log-normalized expression of gene $g$ in class $c$, and $\bar{\mu}_{-c,g} = \frac{1}{|\calT \setminus \calT_c|}\sum_{i \in \calT \setminus \calT_c} x_{ig}$ is the mean across all other labeled training cells. We use $k = 25$ discriminative genes in all experiments.

The \emph{marker score} for cell $i$ and class $c$ is the mean expression over the signature:
\begin{equation}
    s_c(i) = \frac{1}{|\calG_c|}\sum_{g \in \calG_c} x_{ig}.
    \label{eq:marker_score}
\end{equation}

The \emph{marker margin} is:
\begin{equation}
    \Delta_i = s_\rare(i) - s_{\hat{y}_i}(i),
    \label{eq:marker_margin}
\end{equation}
measuring whether candidate cell $i$ expresses rare-class markers more strongly than predicted-class markers.

A validation-set threshold $\tau$ is selected to maximize rare-class F1 subject to a false-rescue rate constraint:
\begin{equation}
    \tau = \argmax_{\tau'} \text{F1}_\rare(\text{val};\tau') \quad \text{subject to} \quad \text{FRR}(\tau') \leq 0.001,
    \label{eq:threshold}
\end{equation}
where $\text{FRR}(\tau')$ is the fraction of non-rare validation cells rescued to the rare class. Only candidates with $\Delta_i > \tau$ are reassigned to label $\rare$.

\subsection{Gated Fusion Variant}
\label{sec:fusion}

As an alternative to the hard gate-and-verify decision, we also evaluate a soft \emph{gated fusion} variant that combines the scANVI softmax probability with a normalized prototype score:
\begin{equation}
    \tilde{p}_\rare(i) = \alpha \cdot P_\phi(\rare \mid \bz_i) + (1-\alpha) \cdot \hat{s}_\rare(i),
    \label{eq:fusion}
\end{equation}
where $\hat{s}_\rare(i)$ is a temperature-softmax-normalized prototype score and $\alpha \in [0,1]$ controls the weight of the scANVI posterior. The prototype gate from \cref{sec:gate} is applied before fusion: if $\sepratio < 1.1$, the fusion score reduces to the scANVI posterior unchanged.

Fusion parameters $(\alpha, T, \theta)$---the scANVI weight, softmax temperature, and rescue probability threshold---are selected on the validation set by a grid search over 18 combinations ($\alpha \in \{0.0, 0.2, 0.4\}$, $T \in \{0.5, 1.0, 2.0\}$, $\theta \in \{0.3, 0.5, 0.7\}$), maximizing rare-class F1 subject to an overall accuracy drop tolerance of 0.5 pp and a false-rescue rate ceiling of 0.5\%.

\subsection{Datasets and Experimental Design}
\label{sec:datasets}

We evaluate scRareRefine on seven datasets representing nine rare cell types (\cref{tab:datasets}). Four datasets use a \emph{batch-held-out} split mode, where entire donors are withheld as the query set (train:validation:test donors split 70:15:15). Three use a \emph{cell-stratified} split mode, where cells are randomly partitioned while preserving class proportions (used when all cells come from a single donor or when batch labels do not correspond to meaningful technical batches). All experiments are repeated across three random seeds (42, 43, 44). The primary evaluation uses $n_\text{rare} = 20$ labeled rare training cells; data-efficiency experiments additionally evaluate $n_\text{rare} \in \{5, 10, 50, \text{all}\}$.

\begin{table}[t]
\centering
\small
\caption{Datasets and rare cell types evaluated in this study. Sep ratio is the mean separability ratio at $n_\text{rare} = 20$ across three seeds. Rare\% is the approximate fraction of rare cells in the full dataset.}
\label{tab:datasets}
\begin{tabular}{lllccc}
\toprule
Dataset & Rare class & Split mode & Sep ratio & Rare\% & Donors \\
\midrule
Immune DC \citep{villani2017single} & ASDC & batch-held-out & 1.50 & $\sim$2\% & 6 \\
Immune DC \citep{villani2017single} & cDC1 & batch-held-out & 1.46 & $\sim$1\% & 6 \\
Human pancreas \citep{baron2016single} & $\varepsilon$-cell & batch-held-out & 1.07 & $\sim$1\% & 4 \\
Human pancreas \citep{baron2016single} & $\gamma$-cell & batch-held-out & 0.92 & $\sim$4\% & 4 \\
Tabula Sapiens Liver \citep{tabulasapiens2022} & NCM & cell-stratified & 1.97 & $\sim$5\% & 1 \\
Tabula Sapiens Pancreas \citep{tabulasapiens2022} & $\beta$-cell & cell-stratified & 0.81 & $\sim$8\% & 1 \\
Tabula Sapiens Spleen \citep{tabulasapiens2022} & ILC & batch-held-out & 1.65 & $\sim$0.2\% & 6 \\
Tabula Sapiens Kidney \citep{tabulasapiens2022} & Endothelial & cell-stratified & 1.69 & $\sim$0.9\% & 1 \\
PBMC COVID atlas \citep{stephenson2021single} & pDC & batch-held-out & 2.03 & $\sim$5\% & 20 \\
\bottomrule
\end{tabular}
\end{table}

Abbreviations: ASDC, AXL$^+$SIGLEC6$^+$ DC precursor; cDC1, conventional dendritic cell type 1; NCM, non-classical monocyte; ILC, innate lymphoid cell; pDC, plasmacytoid dendritic cell.

\paragraph{Baselines.} We compare against: (i) \textbf{scANVI baseline}---unmodified scANVI softmax predictions; (ii) \textbf{kNN ($k$=15)}---$k$-nearest-neighbour classifier fit on scANVI latent embeddings; and (iii) \textbf{LR}---one-versus-rest logistic regression on normalized expression features (CellTypist-equivalent implementation \citep{dominguez2022cross}).

\paragraph{Metrics.} The primary metric is rare-class F1 score on the held-out test set. We also report overall accuracy to confirm that rescue does not degrade majority-class predictions.

\paragraph{Hyperparameters.} scANVI is trained with $n_\text{latent} = 30$, $n_\text{top\_hvg} = 4000$, and default scvi-tools training settings. All other pipeline parameters are fixed as described above and in Supplementary Table~S1.
